\documentstyle[% 


righttag% 


,11pt% 


,amssymb% 


,russian%               remove in Eng. 


,izvru%                 Set izven% in Eng. 


]{amsart} 


 


%       Math definitions 


\newcommand{\Hom}{\operatorname{Hom}} 


\newcommand{\End}{\operatorname{End}} 


\newcommand{\tts}[1]{} 


 


%\hoffset=-15mm%                  For Eng. 


  


\setcounter{page}{77} 


\god{1999} 


\nomer{4~(443)} %                        Remove in Eng. 


%\nomer{Vol.\,43, No.\,4}%               Set in Eng. 


%\str{\pageref{begin}--\pageref{end}}%  Set in Eng. 


\udk{517.983} 


 


\author{\it Ю.В.\,Лазарев, А.В.\,Марченко, И.Б.\,Симоненко} 


% NB:   ^^ remove in Eng. 


\title{О вычислении определителей усеченных одномерных 


континуальных сверток с рациональными символами} 


\thanks{${}^1$\,Под гильбертовым 


пространством мы понимаем линейное 


пространство со скалярным произведением, снабженное 


индуцированной этим произведением нормой и полное 


относительно этой нормы.}


\thanks{$\phantom{.}$}


\thanks{Работа поддержана Российским фондом фундаментальных 


исследований (код проекта 96-01-01195a).} 


 


\date{} 


% \pagestyle{myheadings}%         Set in Eng. 


\pagestyle{plain} %              Remove in Eng. 


\begin{document} 


% \thispagestyle{plain}%          Set in Eng. 


\maketitle 


\label{begin} 


 


 


Имеется много работ, в которых исследовано 


асимптотическое поведение определителя оператора 


усеченной континуальной свертки при увеличении области 


усечения. Прежде всего, это статья [1] и 


монография ([2], с.\,169--172), содержащая достаточно полную 


библиографию по этому вопросу. В данной работе показано, что 


задача вычисления определителя усеченной одномерной 


континуальной свертки с рациональным символом решается в 


квадратурах. 


Вычисления основаны на результатах статьи [3] и 


приведенной ниже теореме. 


 


Пусть $\Bbb N$ --- множество натуральных чисел; $\Bbb R$, $\Bbb C$ --- 


поля 


вещественных и комплексных чисел соответственно; линейное 


пространство --- это пространство, линейное над полем $\Bbb C$; 


$\Hom(B_1,B_2)$, когда $B_1$, $B_2$ --- банаховы пространства, --- 


банахово 


пространство линейных ограниченных операторов, действующих 


из $B_1$ в $B_2$, с естественной нормой; $\End(B)$, когда $B$ --- 


банахово пространство, --- банахова алгебра $\Hom(B,B)$; $I_B$ --- 


единичный оператор из алгебры $\End(B)$; $\frak S_1(B)$, когда $B$ --- 


гильбертово пространство${}^1$, --- 


банахово пространство ядерных 


операторов из $\End(B)$ с естественной для такого пространства 


нормой (напр., [4], с.\,100); $\det(I_B+K)$, когда 


$K\in\frak S_1(B)$, --- определитель оператора $I_B+K$, определяемый как 


$\prod\limits_{j\in\Bbb N}(1+\lambda_j)$, где $\lambda_j$ ($j\in\Bbb N$) 


--- последовательность всех 


собственных чисел оператора $K$ с учетом их кратности. 


Поскольку 


$\sum\limits_{j\in \Bbb N}|\lambda_j|\le\|K\|_{\frak S_1(B)}$ ([5], 


с.\,62), это 


определение корректно и функция $\Phi:\frak S_1(B)\to\Bbb C$, 


определенная 


равенством $\Phi(K)=\det(I_B+K)$, непрерывна. 


 


В данной работе примем следующие соглашения. Под 


функцией будем понимать комплекснозначное отображение;


$[a,b]$, когда $a,b\in\Bbb R$, --- замкнутая выпуклая оболочка точек 


$a,\ b$; 


$(a,b)=[a,b]\setminus\{a,b\}$; $L_2([a,b])$ --- банахово пространство 


суммируемых с квадратом функций, определенных на $[a,b]$, с 


естественной нормой. 


 


Пусть $a,b,c\in\Bbb R$, $c\in(a,b)$; $Y=[a,b]\setminus\{c\}$; $H_y$, 


когда $y\in Y$, --- 


пространство $L_2([c,y])$; $\Pi=[a,b]\times[a,b]$; $k$ --- непрерывная 


на $\Pi$ 


функция; $I_y$,~$K_y$~$(\in\End(H_y))$ ($y\in Y$) --- семейства 


операторов, 


определенные равенствами $I_y f=f$, 


$[K_y(f)](x)=\int\limits^y_c k(x,t)f(t)dt$; 


$Y_0=\{y\mid y\in Y,\text{ оператор }I_y+K_y\text{ обратим}\}$; $k_y$, 


$f_y$ ($y\in Y_0$) --- семейства 


функций, заданных на $[c,y]$ соотношениями $k_y(x)=k(x,y)$, 


$f_y=(I_y+K_y)^{-1}k_y$; $q$ --- функция, определенная на $Y_0$ 


формулой 


$q(y)=f_y(y)$. 


 


Далее, пусть $H=L_2([a,b])$; $K\in\End(H)$ --- оператор, 


определенный равенством $[K(f)](x)=\int\limits^b_a k(x,t)f(t)dt$,


$K\in\frak S_1(H)$; $\Delta$ --- функция, заданная на $Y$ формулой 


$\Delta(y)=\det(I_y+K_y)$; существует такая последовательность 


непрерывных на $\Pi$ функций $r_n$ ($n\in\Bbb N$), 


которая удовлетворяет 


следующим условиям: 


\begin{itemize} 


\item[1)] для каждого $n$ функция $r_n$ представима в виде конечной 


суммы $r_n(x,t)=\sum c_1(x)d_1(t)$, где $c_1$, $d_1$ --- определенные 


на 


сегменте $[a,b]$ непрерывные функции; 


\item[2)] последовательность $r_n$ равномерно стремится к функции 


$k$; 


\item[3)] последовательность $R_n$ ($\in\End(H)$), определенная 


равенством $[R_n(f)](x)=\int\limits^b_a r_n(x,t)f(t)dt$, 


стремится по норме пространства $\frak S_1(H)$ к оператору $K$. 


\end{itemize} 


 


{\bf Теорема.} {\it  


\begin{itemize} 


\item[а)] Mножество $Y_0$ открыто в $Y$. 


\item[б)] Функция $q$ непрерывна. 


\item[в)] Функция $\Delta$ непрерывно дифференцируема на $Y_0$ и при 


$y\in Y_0$ имеет место равенство $\Delta'(y)/\Delta(y)=q(y)$. 


\item[г)] Существует такое $\varepsilon>0$, что 


$(c-\varepsilon,c+\varepsilon)\setminus\{c\}\subset Y_0$, и 


существует предел $\lim\limits_{y\to c}\Delta(y)=1$. 


\end{itemize} 


} 


\medskip 


 


Данная теорема может быть применена для вычисления 


определителей операторов урезанной свертки с рациональным 


символом. Пусть $p\in L_2(-\infty,+\infty)$ --- непрерывная функция; 


преобразование Фурье функции $p$ --- рациональная функция; $I$, $P$ --- 


операторы из алгебры $\End(H)$, определенные равенствами 


$$ 


I f=f,\quad (P f)(x)=\int^b_a p(x-t)f(t)dt. 


$$ 


 


В [3] разработан алгоритм решения уравнения 


$(I+P)f=g$, где $f$ --- искомая, а $g$ --- заданная функции. Используя 


эти результаты и вышеприведенную теорему, можно вычислить 


определитель оператора $I+P$. Например, если $a=0$, 


$b>0$ и $p(x)=\frac{3}{2}\exp(-|x|)$, то $\det(I+P)=(9e^b-e^{-3b})/8$. 


 


\begin{thebibliography}{9} 


 


\bibitem{1} 


Widom H. {\em Szego's limit theorem$:$ the 


higher-dimensional matrix case} // J. Funct. Anal. -- 


1980. -- Т.\,39. -- С.\,182--198. 


\bibitem{2} 


Bottcher A., Silbermann B. {\em Invertibility and 


asymptotics of Toeplitz matrices}. - Berlin: Acad.-Verl., 


1983. -- 200 S. 


\bibitem{3} 


Симоненко И.Б. {\em О некоторых интегро-дифференциальных 


уравнениях типа свертки} // Изв. вузов. 


Математика. -- 1959. -- \No\,2. -- С.\,213--226. 


\bibitem{4} 


Гохберг И.Ц., Крейн М.Г. {\em Теория вольтерровых 


операторов в гильбертовом пространстве и ее приложения}. -- М.: 


Наука, 1967. -- 508 с. 


\bibitem{5} 


Гохберг И.Ц., Крейн М.Г. {\em Введение в теорию линейных 


несамосопряженных операторов}. -- М.: Наука, 1965. -- 448 с. 


 


\end{thebibliography} 


 


\vskip 2cm 


 


\post{\it Ростовский государственный}{\it Поступила} 


\post{\it университет}{19.11.1997} 


 


\label{end} 


\end{document} 


